\chapter{Conclusion}

\section{Summary of Achievements}
This research has successfully established a robust computational framework for traffic sign recognition by synthesizing deep feature extraction from state-of-the-art CNN architectures with high-margin Support Vector Machine classification. The empirical evidence gathered throughout this study validates that transfer learning from ResNet50 and InceptionV3 provides a superior alternative to traditional handcrafted feature engineering. The integration of Global Average Pooling proved to be a critical component in this pipeline, as it effectively projected complex spatial hierarchies onto a 2048-dimensional manifold while preserving essential discriminative class-specific information.

A detailed topological and variance analysis through PCA and UMAP revealed distinct structural properties within these high-dimensional feature spaces. While ResNet50 exhibited higher information density by capturing over half of the global variance within its primary ten components, the multi-scale features of InceptionV3 demonstrated greater expressive power when subjected to non-linear kernel mappings. A pivotal finding of this study is the substantial performance discrepancy observed between linear and non-linear decision boundaries. The transition from linear SVMs to RBF kernels yielded accuracy increments of up to 8.93\%, confirming that traffic sign feature manifolds are intrinsically non-linear and necessitate the application of the RBF kernel trick for optimal class separability. Ultimately, the proposed system achieved a peak classification accuracy of 97.24\%, underscores the effectiveness of combining hierarchical deep features with non-linear maximum-margin classifiers for safety-critical computer vision tasks.

\section{Future Research Directions}
While the current methodology achieves high predictive accuracy, several areas warrant further academic investigation to enhance the system's robustness and applicability. Primary consideration should be given to the implementation of systematic hyperparameter optimization through Bayesian methods, which would likely refine the margin and kernel parameters beyond the standard configurations used here. Furthermore, the model's resilience to adverse environmental conditions could be significantly improved by incorporating domain-specific data augmentation, such as synthetic atmospheric noise or motion blur, during a targeted feature fine-tuning phase. Finally, addressing the trade-off between computational latency and classification precision remains essential for practical deployment. Future research focusing on model quantization or knowledge distillation will be vital to adapting this high-accuracy pipeline for real-time inference on embedded autonomous platforms, ensuring that the system meets the stringent requirements of real-world navigation.
